\documentclass[11pt]{article}
\usepackage[margin=1in]{geometry}
\usepackage{amsmath,amssymb,amsthm,booktabs,longtable,hyperref}
\hypersetup{colorlinks=true,urlcolor=blue,citecolor=blue,linkcolor=blue}
\newtheorem{theorem}{Theorem}
\newtheorem{lemma}{Lemma}
\newcommand{\gs}{\gamma'_{\mathrm{SMB}}}
\newcommand{\gd}{\gamma_{\mathrm{SMB}}}
\title{A seven-vertex obstruction to a minimum-degree bound\\for Staller's fast Maker--Breaker domination game}
\author{AI-assisted research draft\\\small Prepared by the MathDB355140 research lane; no human authorship asserted}
\date{1 November 2026; draft for review, not acceptance or publication}
\begin{document}\emergencystretch=2em\maketitle
\begin{abstract}
We exhibit a simple connected bipartite graph on seven vertices, of minimum degree two, whose optimal Staller-start winning number is four. This exceeds the proposed bound $\lceil 7/2\rceil-2+1=3$ in Bujt\'as and Dokyeesun's Conjecture 7.2(i). A complete finite strategy and an exhaustive three-move obstruction are given as ordinary set arguments. An accompanying Lean dossier proves the exact value and negates the source conjunction. The result concerns clause (i); clause (ii) is separately unresolved here. Literature and priority conclusions are bounded, and this draft awards no discovery credit.
\end{abstract}

\section{Definitions and the exact source question}
Let $G$ be a finite simple graph of order $n>0$, with closed neighborhood
$N[v]=\{v\}\cup\{w:vw\in E(G)\}$ and minimum degree $\delta(G)$.
Dominator and Staller alternately claim one previously unclaimed vertex; passing is not allowed. Write $D$ and $S$ for their disjoint sets. Dominator wins when $D$ is a dominating set, equivalently $D\cap N[v]\ne\varnothing$ for every $v$. Staller wins when $N[v]\subseteq S$ for some $v$. These conditions are incompatible for disjoint occupations. On exhaustion one holds: if no neighborhood lies in $S$, every neighborhood meets its complement $D$.

The number $\gs(G)$ is the least number of \emph{Staller moves} within which she can guarantee a win when she starts, against every legal Dominator strategy, including adversarial delay. It is $\infty$ if she cannot force a win. The unprimed $\gd(G)$ uses Dominator as first player. Total moves are not being counted. These are the definitions in \cite[Sections 1.2--1.3]{BD}, not those of the ordinary domination game.

In both arXiv:2206.12812v1 (26 June 2022) and the published version \cite{BD}, Conjecture 7.2 states, for a graph of order $n$:
\begin{align}
\gs(G)<\infty&\ \Longrightarrow\ \gs(G)\le \lceil n/2\rceil-\delta(G)+1,\tag{i}\label{sourcei}\\
\gd(G)<\infty&\ \Longrightarrow\ \gd(G)\le \lfloor n/2\rfloor-\delta(G)+1.\tag{ii}\label{sourceii}
\end{align}
There is no large-order, connectedness or bipartiteness qualification in that statement. The published body, including the ceiling in (i) and floor in (ii), was inspected in both independent reviews. The MathDB355140 cached record asks for their conjunction; its older bibliographic metadata does not establish the date of this exact conjecture.

\section{The witness and its value}
Let $V=\{0,1,2,3,4,5,6\}$ and
\[
 E=\{01,05,12,16,23,34,45,56\}.
\]
Figure~\ref{fig} displays its three internally disjoint branches between 1 and 5. The graph is simple and connected. Its bipartition is
$\{1,3,5\}\sqcup\{0,2,4,6\}$; every listed edge crosses these parts.
\begin{figure}[ht]\centering
\setlength{\unitlength}{1pt}\begin{picture}(270,120)
\qbezier(10,55)(70.0,80.0)(130,105)
\qbezier(130,105)(190.0,80.0)(250,55)
\qbezier(10,55)(70.0,55.0)(130,55)
\qbezier(130,55)(190.0,55.0)(250,55)
\qbezier(10,55)(40.0,32.5)(70,10)
\qbezier(70,10)(100.0,10.0)(130,10)
\qbezier(130,10)(160.0,10.0)(190,10)
\qbezier(190,10)(220.0,32.5)(250,55)
\put(10,55){\makebox(0,0){\colorbox{white}{1}}}
\put(250,55){\makebox(0,0){\colorbox{white}{5}}}
\put(130,105){\makebox(0,0){\colorbox{white}{0}}}
\put(130,55){\makebox(0,0){\colorbox{white}{6}}}
\put(70,10){\makebox(0,0){\colorbox{white}{2}}}
\put(130,10){\makebox(0,0){\colorbox{white}{3}}}
\put(190,10){\makebox(0,0){\colorbox{white}{4}}}
\end{picture}
\caption{The theta graph with branches $1$--$0$--$5$, $1$--$6$--$5$, and $1$--$2$--$3$--$4$--$5$.}\label{fig}
\end{figure}
\begin{center}\begin{tabular}{ccl}\toprule
$v$ & degree & $N[v]$\\\midrule
0&2&$\{0,1,5\}$\\1&3&$\{0,1,2,6\}$\\2&2&$\{1,2,3\}$\\3&2&$\{2,3,4\}$\\4&2&$\{3,4,5\}$\\5&3&$\{0,4,5,6\}$\\6&2&$\{1,5,6\}$\\\bottomrule
\end{tabular}\end{center}
Thus $\delta(G)=2$.
\begin{theorem}\label{main}
For this graph, $\gs(G)=4$. Consequently Conjecture 7.2(i), and hence the conjunction recorded as MathDB355140, is false.
\end{theorem}

\subsection{A complete upper strategy}
Staller opens at 3. Table~\ref{upper} gives every resulting play under her strategy, stopping as soon as she wins. Columns alternate Staller and Dominator moves. A dash means the game has already ended. The last column gives $v$ with $N[v]\subseteq S$. Repeated histories in the table have identical prescribed Staller moves. At each Dominator prefix the next Dominator column ranges over \emph{every} vertex not previously listed in that history.

These assertions can be read directly: after the opening, $D_1$ ranges over $0,1,2,4,5,6$; after each specified $S_2$, $D_2$ ranges over its four remaining vertices. Only histories $(3,2,5,4,1)$ and $(3,4,1,2,5)$ require $D_3$; each has remaining choices 0 and 6, answered by 6 and 0 respectively. All other histories end at $S_3$. The five possible final Staller sets are $123$, $234$, $345$, $0135$, and $1356$, containing respectively $N[2],N[3],N[4],N[0],N[6]$. Here digit strings denote sets.

Every choice is legal. At any proper prefix, a listed continuation ends with a neighborhood disjoint from the already chosen Dominator vertices, so those vertices cannot yet dominate. Induction backwards through this finite table therefore proves a forced win within four Staller moves. As a rooted tree, counting each S decision, each D branching position, and each terminal leaf once, this certificate has 70 nodes and 26 leaves. No symmetry quotient or unproved symmetry identification is used.

\begin{longtable}{rrrrrrrc}
\caption{Complete upper strategy; all 26 terminal histories.}\label{upper}\\
\toprule $S_1$&$D_1$&$S_2$&$D_2$&$S_3$&$D_3$&$S_4$&$v$\\\midrule\endfirsthead
\toprule $S_1$&$D_1$&$S_2$&$D_2$&$S_3$&$D_3$&$S_4$&$v$\\\midrule\endhead
3 & 0 & 2 & 1 & 4 & -- & -- & 3 \\
3 & 0 & 2 & 4 & 1 & -- & -- & 2 \\
3 & 0 & 2 & 5 & 1 & -- & -- & 2 \\
3 & 0 & 2 & 6 & 1 & -- & -- & 2 \\
3 & 1 & 4 & 0 & 2 & -- & -- & 3 \\
3 & 1 & 4 & 2 & 5 & -- & -- & 4 \\
3 & 1 & 4 & 5 & 2 & -- & -- & 3 \\
3 & 1 & 4 & 6 & 2 & -- & -- & 3 \\
3 & 2 & 5 & 0 & 4 & -- & -- & 4 \\
3 & 2 & 5 & 1 & 4 & -- & -- & 4 \\
3 & 2 & 5 & 4 & 1 & 0 & 6 & 6 \\
3 & 2 & 5 & 4 & 1 & 6 & 0 & 0 \\
3 & 2 & 5 & 6 & 4 & -- & -- & 4 \\
3 & 4 & 1 & 0 & 2 & -- & -- & 2 \\
3 & 4 & 1 & 2 & 5 & 0 & 6 & 6 \\
3 & 4 & 1 & 2 & 5 & 6 & 0 & 0 \\
3 & 4 & 1 & 5 & 2 & -- & -- & 2 \\
3 & 4 & 1 & 6 & 2 & -- & -- & 2 \\
3 & 5 & 2 & 0 & 1 & -- & -- & 2 \\
3 & 5 & 2 & 1 & 4 & -- & -- & 3 \\
3 & 5 & 2 & 4 & 1 & -- & -- & 2 \\
3 & 5 & 2 & 6 & 1 & -- & -- & 2 \\
3 & 6 & 2 & 0 & 1 & -- & -- & 2 \\
3 & 6 & 2 & 1 & 4 & -- & -- & 3 \\
3 & 6 & 2 & 4 & 1 & -- & -- & 2 \\
3 & 6 & 2 & 5 & 1 & -- & -- & 2 \\
\bottomrule\end{longtable}

\subsection{An exhaustive lower strategy}
Only five closed neighborhoods have size three:
\[
\mathcal T=\{015,123,234,345,156\}.
\]
The other two have size four. A win before three Staller moves is impossible. After two Staller moves $a,b$, let
\[
 C(a,b)=\{c\notin\{a,b\}:\{a,b,c\}\in\mathcal T\}.
\]
Dominator responds to the opening $a$ with $d$, and to the second move $b$ with $e$, as in Table~\ref{lower}. Its first two columns are constant within each opening block. The five possible $b$ in each block are exactly all legal second Staller choices. Each displayed $e$ is still unclaimed. Let $F=V\setminus\{a,d,b,e\}$ be the three remaining vertices. The last two columns explicitly give $F$ and $C(a,b)$; they are disjoint in every row.

For each of these 35 positions and each of the three legal third Staller choices $c\in F$, her three vertices therefore contain \emph{none of the seven closed neighborhoods}. Indeed, containing a size-three neighborhood would mean $c\in C(a,b)$, and a size-four neighborhood cannot lie in a three-element set. This covers all 105 completions. If Dominator has already won, play simply stops and Staller still has not won. The table is thus a legal obstruction against every Staller opening and every second choice, proving $\gs(G)>3$. Combined with the upper strategy it proves Theorem~\ref{main}.

{\renewcommand{\arraystretch}{0.93}\begin{longtable}{rrrrcc}
\caption{Lower strategy and complete completion-set check. Digit strings denote sets.}\label{lower}\\
\toprule $a=S_1$&$d=D_1$&$b=S_2$&$e=D_2$&$F$&$C(a,b)$\\\midrule\endfirsthead
\toprule $a$&$d$&$b$&$e$&$F$&$C(a,b)$\\\midrule\endhead
0 & 1 & 2 & 3 & 456 & $\varnothing$ \\
0 & 1 & 3 & 2 & 456 & $\varnothing$ \\
0 & 1 & 4 & 2 & 356 & $\varnothing$ \\
0 & 1 & 5 & 2 & 346 & 1 \\
0 & 1 & 6 & 2 & 345 & $\varnothing$ \\
1 & 0 & 2 & 3 & 456 & 3 \\
1 & 0 & 3 & 2 & 456 & 2 \\
1 & 0 & 4 & 2 & 356 & $\varnothing$ \\
1 & 0 & 5 & 6 & 234 & 06 \\
1 & 0 & 6 & 5 & 234 & 5 \\
2 & 1 & 0 & 3 & 456 & $\varnothing$ \\
2 & 1 & 3 & 4 & 056 & 14 \\
2 & 1 & 4 & 3 & 056 & 3 \\
2 & 1 & 5 & 0 & 346 & $\varnothing$ \\
2 & 1 & 6 & 0 & 345 & $\varnothing$ \\
3 & 2 & 0 & 1 & 456 & $\varnothing$ \\
3 & 2 & 1 & 0 & 456 & 2 \\
3 & 2 & 4 & 5 & 016 & 25 \\
3 & 2 & 5 & 4 & 016 & 4 \\
3 & 2 & 6 & 0 & 145 & $\varnothing$ \\
4 & 2 & 0 & 1 & 356 & $\varnothing$ \\
4 & 2 & 1 & 0 & 356 & $\varnothing$ \\
4 & 2 & 3 & 5 & 016 & 25 \\
4 & 2 & 5 & 3 & 016 & 3 \\
4 & 2 & 6 & 0 & 135 & $\varnothing$ \\
5 & 0 & 1 & 6 & 234 & 06 \\
5 & 0 & 2 & 1 & 346 & $\varnothing$ \\
5 & 0 & 3 & 4 & 126 & 4 \\
5 & 0 & 4 & 3 & 126 & 3 \\
5 & 0 & 6 & 1 & 234 & 1 \\
6 & 0 & 1 & 5 & 234 & 5 \\
6 & 0 & 2 & 1 & 345 & $\varnothing$ \\
6 & 0 & 3 & 1 & 245 & $\varnothing$ \\
6 & 0 & 4 & 1 & 235 & $\varnothing$ \\
6 & 0 & 5 & 1 & 234 & 1 \\
\bottomrule\end{longtable}}

\subsection{The source consequence and its limits}
The value is finite and
\[
4=\gs(G)>\lceil7/2\rceil-2+1=3.
\]
This is a counterexample to (i) under its literal hypotheses. Clause (ii) is not separately resolved. Independent minimax diagnostics give $\gd(G)=\infty$ here; this diagnostic does not refute its conditional bound and is not a second claim of novelty.

Proposition 7.1 of \cite{BD} gives the unrestricted ceiling bound for a finite Staller-start value; its displayed sharp constructions have minimum degree one. The source specifically motivates Conjecture 7.2 by the lack of sharp examples of degree at least two. This witness attains the unrestricted ceiling at degree two, so the proposed one-move degree improvement fails at that degree. This is the substantive consequence, independently of runtime, certificate shortness or an open database label. No minimum counterexample order, infinite family, asymptotic obstruction, new universal bound or second discovery is asserted.

\section{Comparison with inspected primary results}
The following scope comparisons follow the first independent source/significance review, whose full report and exact primary receipts accompany the dossier. They record inspected theorem bodies, not an exhaustive literature survey.

In \cite{BD}, Lemma 3.1 and Theorem 3.2 give realization constructions: the seven-vertex $F'_2$ has Staller value three, so it does not give this value four. The path/cycle results and tadpole Theorem 6.2 do not cover a leafless theta with two independent cycles. One-sided cutting or edge-comparison bounds alone do not supply this lower bound. In \cite{trees}, Theorem 3.5 and Proposition 3.7/Corollary 3.8 concern trees or degree-preserving tree substructures; required ambient leaves cannot occur here. Subdivided-star formulas do not yield this witness.

The union/join outcome theorems of \cite{original} concern outcomes rather than optimal Staller speed; the witness and its complement are connected. The residual-matching and tree/cycle formulas of \cite{numbers} primarily concern Dominator's fast-winning parameter. The pairing and twins tools of \cite{partition} can simplify outcomes but do not preserve optimal Staller timing; its outerplanar characterization does not apply because contracting the long branch gives a $K_{2,3}$ minor. Biased-game degree thresholds in \cite{biased} do not give this Staller-start value. The predominated critical tree/cactus and odd-branch constructions of \cite{critical} do not supply this even-branch noncactus timing result. The minimum-order outcome threshold $\beta'(2)=7$ in \cite{minimum} is not a speed formula; its universal-vertex construction is not this maximum-degree-three graph. Generic short-Staller algorithms and complexity results in \cite{complexity} are not an explicit earlier evaluation or refutation.

No full prior subsumer was located in this bounded comparison. The independent judgment is \emph{qualified PASS} for novelty/significance, not global priority certification. Specialist/citing bodies and thesis details not fully inspected, inaccessible live MathDB/Palomar solution state (inherited HTTP 403), deleted histories and most registry proof bodies remain outside clearance. The 1 November 2026 review refreshed the following registry heads:\par OpenAI/math:
\texttt{4eaee87548250b2682af}\allowbreak\texttt{fda7d05d947c05554bae}
\par Omega:
\texttt{a87bdf39e4517497a977}\allowbreak\texttt{3c0c724243bc700c5bb1}.
Complete specified path inventories and selected near-name bodies were screened; absence of a path hit is not an all-proof or all-history search. No new network refresh was undertaken for this manuscript. Dates/content, rather than current searchability alone, govern chronology. Exact receipts and remaining limits are preserved with the reports.

\section{Lean semantics, theorem scope, and reproduction}
The unchanged \texttt{Game.lean} models positions by Staller's set $s$ and the available set $a$. On legal initial play, $D=V\setminus(s\cup a)$, and the three sets partition the board. A Staller move inserts its chosen available vertex into $s$ and erases it from $a$; a Dominator move only erases its chosen vertex. This preserves legal occupations and disallows repeated vertices or passes.

The inductive proposition \texttt{Wins G s a turn b} represents a finite strategy guaranteeing a Staller win within budget $b$. Its \texttt{done} constructor requires an owned closed neighborhood. Its Staller constructor chooses one legal vertex and spends one budget unit. Its Dominator constructor requires a nonempty board and a continuation for every legal choice, spending no Staller budget. This nonemptiness prevents a vacuous win on an exhausted board.

Theorems \texttt{terminal\_incompatibility} and \texttt{Wins.not\_after\_dominator} exclude every successful continuation after Dominator owns a dominating set. They provide external semantic justification for omitting explicit early Dominator termination; they are \emph{not dependencies} of the computational endpoint. Finiteness of play (at most seven total moves here) identifies finite strategy trees with forced wins in the actual game.

\texttt{check\_iff} proves the recursive Boolean checker equivalent to \texttt{Wins} for arbitrary finite graphs, positions, turns and budgets, by induction on the available-set cardinality. The zero-budget, exhausted-board and done cases are included. \texttt{four\_check} and \texttt{three\_check} use kernel decision, and \texttt{Wins.mono} makes budgets mean ``at most.'' The noncomputable \texttt{winningNumber} is the least successful budget in $\mathbb N\cup\{\infty\}$, or top when none exists.

The separately attributed, unchanged independent reviewer supplement \texttt{Bridge.lean} proves
\[
\mathrm{winningNumber}(G,t)\le B
\quad\Longleftrightarrow\quad
\mathrm{Wins}(G,\varnothing,V,t,B)
\]
for every finite graph, turn and natural budget, as \texttt{Review.number\_le\_iff}. It also supplies empty-board and singleton edge controls. This supplement is distinct from the original 63 declarations. The draft replay audits all 63 main declarations and the supplement's sole retained declaration, \texttt{Review.number\_le\_iff}, including any generated declarations in either module. Anonymous edge-control examples are checked during compilation and are not retained as declarations.

The endpoints, without additional assumptions, are
\begin{quote}\ttfamily
MathDB355140.winning\_number\_four\\
MathDB355140.source\_refutation.
\end{quote}
The first states equality to four in \texttt{WithTop Nat}. The second negates \texttt{SourceClaim}, the universal conjunction using successful budgets and additive bounds $b+\delta\le\lceil n/2\rceil+1$ (and floor). Via the budget bridge these are the source's integer inequalities; the additive form avoids truncated subtraction. For the witness they require $b\le3$, contradicted by the lower-budget exclusion and monotonicity. The auxiliary natural-subtraction \texttt{BudgetClaim} is not needed for this source endpoint. No arbitrary-label relabeling theorem is needed to exhibit a graph on \texttt{Fin 7}.

\subsection{Reproduction and trust boundary}
Extract \texttt{355140-proof-source.zip} into a fresh directory, then run \texttt{python3 replay.py}. The source-only packet contains both unchanged modules, an all-declaration auditor, ordinary verifier/tree, pins and instructions; no prebuilt local oleans are included. The managed replay builds both modules with \texttt{-j1 -DwarningAsError=true}, then runs the auditor and module \texttt{leanchecker} checks for both. Full explicit types, available values/proof terms and individual transitive axiom reports are retained. Only \texttt{propext}, \texttt{Classical.choice}, and \texttt{Quot.sound} are allowed. Safety flags and local axiom declarations are checked, including generated declarations. The generated name \texttt{check.\_unsafe\_rec} has unsafe flag false. No sorry, admit, new axiom, unsafe declaration, native decision or \texttt{ofReduceBool} is used.

Lean is pinned to 4.34.1, commit
\texttt{5045d0056413266e57c6}\allowbreak\texttt{25dcd7c365b10e377c52}; mathlib to
\texttt{d13f23b723b8a846827a}\allowbreak\texttt{245b89c10fc7d3f11612}.
The default binary and external cache paths, binary/manifests/direct-import cache hashes and complete commands are in the packet. Each invocation has at most two CPUs, a 6 GiB sampled RSS cap, 12 GiB address-space cap, 120 CPU seconds, 140 wall seconds, and \texttt{MALLOC\_ARENA\_MAX=2}. Sampling every 25 ms can miss transient peaks; this is not an aggregate cgroup guarantee. Invocations run sequentially in a sanitized environment. Negative controls reject a false three-move win, a missing module and a corrupt module; the supplement's positive edge controls compile strictly.

The external prebuilt upstream oleans remain trusted. Cache/source hashes identify inputs but do not prove source-to-olean correspondence. No standalone dependencies or full upstream source rebuild is claimed. Module leanchecker uses the same Lean kernel implementation, not an independent foundation. Ordinary correctness rests also on the explicit tables and their finite set proof, not solely compiler success. Running \texttt{python3 ordinary-proof.py} independently rechecks the tree and the 105 completions with finite resource limits.

\section{Draft status and evidence}
Both independent reviews and root integrations pass ordinary mathematics, inspected source versions, fidelity and strict Lean gates; novelty/significance retain their qualified scope. The root corrected exactly one ambiguous ordinary-proof sentence to ``contains none of the seven closed neighborhoods.'' This manuscript uses that corrected revision. Game bytes and the original source ZIP remain frozen. The dossier verifies all 213 frozen bindings before and after preparation, retains replay logs, and references original evidence by path and hash without recursively copying historical trees. Manuscript review, final root replay and acceptance are separate. Credit here is zero; the campaign remains 3/20 beyond 39. No publication, upload or external communication is requested or performed.

\begin{thebibliography}{10}\raggedright
\bibitem{BD} C. Bujt\'as and P. Dokyeesun, \emph{Fast winning strategies for Staller in the Maker--Breaker domination game}, Discrete Applied Mathematics 344 (2024), 10--22. DOI: \href{https://doi.org/10.1016/j.dam.2023.11.015}{10.1016/j.dam.2023.11.015}. Preprint arXiv:2206.12812v1, 26 June 2022.
\bibitem{trees} C. Bujt\'as, P. Dokyeesun and S. Klav\v zar, \emph{Maker--Breaker domination game on trees when Staller wins}, Discrete Mathematics \& Theoretical Computer Science 25:2, no. 12 (2023); arXiv:2212.06530v3, 25 August 2023.
\bibitem{original} E. Duch\^ene, V. Gledel, A. Parreau and G. Renault, \emph{Maker--Breaker domination game}, arXiv:1807.09479v2, 18 September 2018.
\bibitem{numbers} V. Gledel, V. Ir\v si\v c and S. Klav\v zar, \emph{Maker--Breaker domination number}, arXiv:1810.04397v2, 13 March 2019.
\bibitem{partition} G. Bagan, E. Duch\^ene, V. Gledel, T. Lehtil\"a and A. Parreau, \emph{Partition strategies for the Maker--Breaker domination game}, arXiv:2406.15165v1, 21 June 2024.
\bibitem{biased} A. D. Bagdas, D. Clemens, F. Hamann and Y. Mogge, \emph{Biased domination games}, arXiv:2408.00529v1, 1 August 2024.
\bibitem{critical} C. Bujt\'as, P. Dokyeesun, S. Klav\v zar and M. Stojakovi\'c, \emph{Criticality for Maker--Breaker domination games with predomination}, arXiv:2503.11907v1, 14 March 2025.
\bibitem{minimum} J. F\"uhrer, G. Grasegger, P. Hametner and O. Roche-Newton, \emph{The minimum degree question for the Maker Breaker Domination Game}, arXiv:2606.04824v1, 3 June 2026.
\bibitem{complexity} G. Bagan, M. Hilaire, N. Oijid and A. Parreau, \emph{On the parameterized complexity of the Maker-Breaker domination game}, arXiv:2601.08562v1, 13 January 2026.
\end{thebibliography}
\end{document}
